\section*{الفصل \ref{ch:g5:natural-hazards} --- الأخطار الطبيعية}

\begin{solution}{exo:g5:natural-hazards:1}
(أ) \omterm{def:g5:natural-hazards:hazard}{خطر طبيعي} فقط: لا أحد يتأذى. (ب) \omterm{def:g5:natural-hazards:disaster}{كارثة طبيعية}: البيوت
تتضرر. (ج) \omterm{def:g5:natural-hazards:hazard}{خطر طبيعي} فقط. (د) \omterm{def:g5:natural-hazards:disaster}{كارثة طبيعية}.
\end{solution}

\begin{solution}{exo:g5:natural-hazards:2}
من الهواء: موجة حر، \omterm{def:g5:natural-hazards:storm}{عاصفة}، \omterm{def:g5:natural-hazards:drought}{جفاف}. من الماء: \omterm{def:g5:water-shapes-the-land:flood}{فيضان}. من
الأرض: \omterm{def:g5:natural-hazards:landslide}{انزلاق أرضي}، \omterm{def:g4:volcanoes-earthquakes-first-look:eruption}{ثوران} بركاني، \omterm{def:g4:volcanoes-earthquakes-first-look:earthquake}{زلزال}.
\end{solution}

\begin{solution}{exo:g5:natural-hazards:3}
\omterm{def:g5:natural-hazards:drought}{الجفاف} فترة طويلة، من أشهر أو سنين، يسقط فيها من \omterm{def:g4:clouds-rain-snow:precipitation}{المطر} أقل بكثير من
المعتاد. وقد يترك المزارعين بلا حصاد، لأن المحاصيل تذبل
وتموت، وقد يترك المدن في نقص من ماء الشرب، لأن \omterm{def:g1:landscapes-around-us:water}{الأنهار}
\omterm{def:g1:landscapes-around-us:water}{والبحيرات} تنكمش.
\end{solution}

\begin{solution}{exo:g5:natural-hazards:4}
مثلًا: أدخِلوا أو اربطوا كل ما قد تحمله الريح بعيدًا
(كراسي الحديقة، الألعاب، صناديق القمامة)؛ وتأكدوا من المصابيح اليدوية وجهّزوا بطاريات احتياطية
وماءً؛ وابقوا في الداخل، بعيدًا عن النوافذ، أثناء \omterm{def:g5:natural-hazards:storm}{العاصفة}؛ وواصلوا
الاستماع إلى المذياع لمعرفة الأخبار.
\end{solution}

\begin{solution}{exo:g5:natural-hazards:5}
(أ) هاريكين. (ب) تيفون. (ج) إعصار مداري.
\end{solution}

\begin{solution}{exo:g5:natural-hazards:6}
أخضر، أصفر، برتقالي، أحمر. عند المستوى الأحمر، يكون الخطر كبيرًا: تصرّف
فورًا واتبع تعليمات المسؤولين.
\end{solution}

\begin{solution}{exo:g5:natural-hazards:7}
انفصل جزء من سفح التل إلى اليمين وانزلق، تاركًا ندبة
عارية. وانتشر وحله عبر قاع \omterm{def:g1:landscapes-around-us:landform}{الوادي} كله، فغطى الطريق. كما
سد \omterm{def:g1:landscapes-around-us:water}{النهر} كأنه سد: فماء \omterm{def:g1:landscapes-around-us:water}{النهر}، الذي لم يعد يستطيع أن
يجري بعيدًا، تكدس خلفه وغمر الغابة في أعلى المجرى.
\end{solution}

\begin{solution}{exo:g5:natural-hazards:8}
لينا محقة. لا نستطيع أن نوقف \omterm{def:g4:clouds-rain-snow:precipitation}{المطر} أو \omterm{def:g1:landscapes-around-us:water}{النهر}، لكننا نستطيع أن نبعد الناس
عن الأذى. مثلًا: أن نبني البيوت على أرض أعلى بدلًا من
السهل الفيضي؛ وأن نحذر الناس في الوقت المناسب بتحذيرات \omterm{def:g5:water-shapes-the-land:flood}{الفيضان}، لكي ينتقلوا
إلى أرض أعلى؛ وأن نعلّم الجميع ألا يعبروا أبدًا مياه \omterm{def:g5:water-shapes-the-land:flood}{الفيضان} مشيًا أو
بالسيارة.
\end{solution}

\begin{solution}{exo:g5:natural-hazards:9}
(أ) $4 \times 60 = 240$ مليمترًا. (ب) ربع 60 هو 15؛ وفي أربعة
أشهر، $4 \times 15 = 60$ مليمترًا. (ج) نقص $240 - 60 = 180$ مليمترًا:
أي ثلاثة أرباع \omterm{def:g4:clouds-rain-snow:precipitation}{المطر} المعتاد.
\end{solution}

\begin{solution}{exo:g5:natural-hazards:10}
(أ) $1\,200 \div 4 = 300$ مسكن. (ب) ثلث 300 هو 100؛ إذن يبقى
$300 - 100 = 200$ مسكن في \omterm{def:g5:water-shapes-the-land:flood}{السهل الفيضي}.
\end{solution}

\begin{solution}{exo:g5:natural-hazards:11}
أسباب ممكنة: القرية الأولى حُذّرت في الوقت المناسب واستعد
أهلها (أغلقوا المصاريع، وأدخلوا الأشياء)، والثانية لم تُحذَّر؛ وبيوت
الأولى مبنية بمتانة، وبيوت الثانية هشة؛ والقرية
الثانية تقع على \omterm{def:g1:landscapes-around-us:water}{ساحل} منخفض أو في \omterm{def:g5:water-shapes-the-land:flood}{سهل فيضي}، حيث جاءت \omterm{def:g5:natural-hazards:storm}{العاصفة} أيضًا
\omterm{def:g5:water-shapes-the-land:flood}{بفيضان}، بينما تقوم الأولى على أرض أعلى؛ والأشجار القريبة من
البيوت سقطت عليها في القرية الثانية. فالخطر نفسه لا يُحدث
الأذى نفسه في كل مكان: إنه يتوقف على مكان عيش الناس وطريقته، وعلى
مدى استعدادهم.
\end{solution}
